% =====================================================================
%  capitulos/00b_nomenclatura.tex  --  Lista de simbolos y abreviaturas
%  Material preliminar. Se incluye desde main.tex despues de \listoftables
%  y antes de \pagenumbering{arabic}. Tabla manual (longtable) para no
%  depender del paquete nomencl.
% =====================================================================

\chapter*{Nomenclatura}
\addcontentsline{toc}{chapter}{Nomenclatura}
\markboth{Nomenclatura}{Nomenclatura}

Se reúnen los símbolos y las abreviaturas empleados en el documento. Las
unidades se indican en el Sistema Internacional como referencia dimensional;
en los casos de validación se emplean además sistemas técnicos coherentes
(por ejemplo, kgf y cm, con tensiones en kgf/cm$^2$), según se aclara en cada problema.
Los vectores y matrices se escriben en negrita: así, $t$ es el espesor y $\bm{t}$ el
vector de tracciones. \textbf{Advertencia de notación:} $E$, sin negrita, es el módulo
de elasticidad, mientras que $\bm{E}$, en negrita, es la matriz de extrapolación de
tensiones a los nodos; son magnitudes distintas y la negrita es lo único que las separa.

Los símbolos no se listan en orden alfabético: se agrupan por tema, y dentro de cada
tema en el orden en que aparecen en el desarrollo del documento.

\section*{Símbolos}

\renewcommand{\arraystretch}{1.2}
\begin{longtable}{@{}>{\raggedright\arraybackslash}p{3.0cm}p{7.8cm}p{3.0cm}@{}}
\toprule
\textbf{Símbolo} & \textbf{Significado} & \textbf{Unidad} \\
\midrule
\endfirsthead
\toprule
\textbf{Símbolo} & \textbf{Significado} & \textbf{Unidad} \\
\midrule
\endhead
\bottomrule
\endfoot
\multicolumn{3}{@{}l}{\textit{Dominio, material y campos}} \\
\addlinespace
$\Omega,\ \Gamma$ & Dominio del problema y su frontera & --- \\
$\Gamma_u,\ \Gamma_t$ & Frontera con desplazamientos / tracciones prescritos & --- \\
$(x,y)$ & Coordenadas físicas (cartesianas) & m \\
$(\xi,\eta)$ & Coordenadas naturales del elemento de referencia & adimensional \\
$E$ & Módulo de elasticidad (de Young) & Pa \\
$\nu$ & Coeficiente de Poisson & adimensional \\
$\rho$ & Densidad del material & kg/m$^3$ \\
$t$ & Espesor del elemento & m \\
$h$ & Tamaño característico de malla & m \\
$p$ & Orden polinómico del elemento ($1$ en Q4, $2$ en Q9) & adimensional \\
$\bm{u}$ & Vector de desplazamientos nodales & m \\
$\bm{u}_e$ & Desplazamientos de los nodos de un elemento & m \\
$u,\ v$ & Componentes de desplazamiento en $x$ e $y$ & m \\
$\delta\bm{u}$, $\delta\bm{\varepsilon}$ & Desplazamientos y deformaciones virtuales del principio de los trabajos virtuales & m; adimensional \\
$\bm{b}$ & Fuerza por unidad de volumen & N/m$^3$ \\
$\bm{t}$ & Vector de tracciones de borde & N/m$^2$ \\
$\bm{\sigma}$ & Vector de tensiones $[\sigma_x,\sigma_y,\tau_{xy}]^T$ (Voigt) & Pa \\
$\sigma_x,\ \sigma_y$ & Tensiones normales & Pa \\
$\tau_{xy}$ & Tensión cortante & Pa \\
$\sigma_z$ & Tensión normal fuera del plano ($0$ en tensión plana; $\nu(\sigma_x+\sigma_y)$ en deformación plana) & Pa \\
$\sigma_1,\ \sigma_2$ & Tensiones principales & Pa \\
$\sigma_{\mathrm{VM}}$ & Tensión equivalente de von Mises & Pa \\
$\bm{\varepsilon}$ & Vector de deformaciones $[\varepsilon_x,\varepsilon_y,\gamma_{xy}]^T$ & adimensional \\
$\varepsilon_x,\ \varepsilon_y$ & Deformaciones normales & adimensional \\
$\gamma_{xy}$ & Deformación angular (de corte) & adimensional \\
\addlinespace
\multicolumn{3}{@{}l}{\textit{Interpolación y geometría del elemento}} \\
\addlinespace
$e,\ i,\ j,\ g$ & Índices: $e$ recorre los elementos, $i$ y $j$ los nodos y $g$ los puntos de cuadratura & --- \\
$n$ & Número de nodos del elemento ($4$ en Q4, $9$ en Q9) & adimensional \\
$N_i$ & Función de forma asociada al nodo $i$ & adimensional \\
$\bm{N}$ & Matriz de funciones de forma & adimensional \\
$\bm{X}_e$ & Coordenadas de los nodos del elemento & m \\
$\bm{J}$ & Matriz Jacobiana del mapeo isoparamétrico & m \\
$\det\bm{J}$ & Determinante del Jacobiano & m$^2$ \\
$\bm{B}$ & Matriz de deformación-desplazamiento & m$^{-1}$ \\
$\bm{D}$ & Matriz constitutiva (elástica); $\bm{D}_{\text{TP}}$ en tensión plana y $\bm{D}_{\text{DP}}$ en deformación plana & Pa \\
\addlinespace
\multicolumn{3}{@{}l}{\textit{Cuadratura, rigidez y ensamblaje}} \\
\addlinespace
$w_g$ & Peso de la cuadratura de Gauss en el punto $g$ & adimensional \\
$n_g$ & Número de puntos de Gauss del elemento & adimensional \\
$\bm{k}_e$ & Matriz de rigidez elemental & N/m \\
$\Ensamblaje$ & Operador de ensamblaje de las matrices elementales en el sistema global & --- \\
$\bm{K}$ & Matriz de rigidez global & N/m \\
$\bm{F}$ & Vector de cargas global & N \\
$\bm{F}_j$ & Carga nodal aplicada en el nodo $j$ & N \\
$\bm{f}_e^{\text{vol}}$ & Cargas volumétricas consistentes del elemento & N \\
$\bm{F}_\Gamma$ & Fuerzas nodales equivalentes de una carga de superficie & N \\
$L$ & Longitud de la arista cargada del elemento & m \\
$q_1,\ q_2$ & Valores de la carga distribuida en los extremos de esa arista & N/m \\
$\bm{K}_{ff}$, $\bm{K}_{fr}$, $\bm{K}_{rf}$, $\bm{K}_{rr}$ & Bloques de $\bm{K}$ tras partir los GDL en libres ($f$) y restringidos ($r$) & N/m \\
$\bm{u}_f,\ \bm{u}_r$ & Desplazamientos de los GDL libres y de los restringidos & m \\
$\bm{R}$ & Vector de reacciones en los apoyos & N \\
$\kappa_2$ & Número de condición de la matriz de rigidez, indicador de cuán mal planteado está el sistema & adimensional \\
\addlinespace
\multicolumn{3}{@{}l}{\textit{Recuperación de tensiones}} \\
\addlinespace
$\bm{M}$ & Matriz de las funciones de forma evaluadas en los puntos de Gauss, cuya inversa es $\bm{E}$ & adimensional \\
$\bm{E}$ & Matriz de extrapolación de tensiones a los nodos. No confundir con el módulo de elasticidad $E$, que se escribe sin negrita & adimensional \\
$n_e$ & Número de elementos que comparten un nodo, en el promediado nodal & adimensional \\
\addlinespace
\multicolumn{3}{@{}l}{\textit{Calidad geométrica de la malla}} \\
\addlinespace
$q_{SJ}$ & Jacobiano escalado (métrica de calidad, $\in[-1,1]$) & adimensional \\
$q_C$ & Compacidad o \emph{stretch} (métrica de calidad, $\in[0,1]$) & adimensional \\
$L_{\min}$ & Longitud de la arista más corta del elemento & m \\
$D_{\max}$ & Longitud de la diagonal más larga del elemento & m \\
$R_J$ & Razón entre el menor y el mayor determinante del Jacobiano del elemento & adimensional \\
$\bm{r}_1,\dots,\bm{r}_4$ & Coordenadas de las cuatro esquinas del cuadrilátero & m \\
$\bm{X}_1$, $\bm{X}_2$, $\bm{X}_3$ & Vectores de la descomposición bilineal del cuadrilátero & m \\
$AR$ & Relación de aspecto del elemento & adimensional \\
$T_R$ & Estrechamiento o \emph{taper} del elemento & adimensional \\
\addlinespace
\multicolumn{3}{@{}l}{\textit{Verificación y validación}} \\
\addlinespace
$\bm{u}_M$ & Campo de desplazamientos exacto manufacturado (MMS) & m \\
$\bm{u}_h$ & Solución numérica sobre una malla de tamaño característico $h$ & m \\
$\bm{\sigma}^*_h$ & Campo de tensiones recuperado: extrapolado a los nodos y promediado & Pa \\
$N$ & Número de subdivisiones por lado en las mallas estructuradas de refinamiento ($N\times N$ elementos). No confundir con la función de forma $N_i$ & adimensional \\
$\|\cdot\|_{L^2}$ & Norma $L^2$ del error en desplazamiento & unidad de la magnitud \\
$|\cdot|_{H^1}$ & Seminorma $H^1$ del error (gradiente) & unidad de la magnitud / m \\
\end{longtable}

\section*{Abreviaturas y siglas}

\begin{longtable}{@{}p{2.4cm}p{11.4cm}@{}}
\toprule
\textbf{Sigla} & \textbf{Significado} \\
\midrule
\endfirsthead
\toprule
\textbf{Sigla} & \textbf{Significado} \\
\midrule
\endhead
\bottomrule
\endfoot
MEF & Método de los Elementos Finitos \\
GDL & Grado(s) de libertad \\
TP & Tensión plana \\
DP & Deformación plana \\
Q4 & Elemento cuadrilátero isoparamétrico de 4 nodos (bilineal) \\
Q9 & Elemento cuadrilátero isoparamétrico de 9 nodos (bicuadrático) \\
MMS & Método de las Soluciones Manufacturadas \\
V\&V & Verificación y Validación \\
LU & Factorización LU (descomposición triangular inferior-superior) \\
COO & Formato disperso de coordenadas (\emph{coordinate list}) \\
CSR & Formato disperso de filas comprimidas (\emph{compressed sparse row}) \\
SRI & Integración reducida selectiva (\emph{selective reduced integration}) \\
MVC & Patrón de arquitectura Modelo-Vista-Controlador \\
CAD & Diseño asistido por computadora (\emph{computer-aided design}) \\
DXF & Formato de intercambio de dibujo CAD (\emph{drawing exchange format}) \\
CSV & Valores separados por comas (\emph{comma-separated values}) \\
PDF & Formato de documento portátil (\emph{portable document format}) \\
ICMJE & Comité Internacional de Editores de Revistas Médicas (\emph{International Committee of Medical Journal Editors}), responsable de la norma Vancouver \\
NLM & Biblioteca Nacional de Medicina de los Estados Unidos (\emph{National Library of Medicine}) \\
PI & Pregunta de investigación (PI-1 a PI-3) \\
FEA, FEM & Códigos de los ítems del consenso de expertos que emplea la tabla de cobertura de contenido; no se usan como siglas en la prosa \\
\end{longtable}
\renewcommand{\arraystretch}{1.0}

% Las dos longtable anteriores no llevan \caption, pero igualmente incrementan
% el contador `table`. Se reinicia para que la numeracion de tablas del cuerpo
% no quede desplazada (p. ej. que la primera tabla de la Introduccion no salga
% como "Tabla 3").
\setcounter{table}{0}
